% Rapport du projet 1 -- Robot de reconfort (structure imposee : enonce section 11)
% 3 pages MAXIMUM, quatre points et rien d'autre : pas de page de titre, pas de
% table des matieres, pas d'introduction ni de conclusion.
% Les commentaires « % → » rappellent ce que l'enonce attend, et les « % données : »
% sont les chiffres mesures sur notre code ; a supprimer une fois redige.

\documentclass[11pt,a4paper]{article}
\usepackage[utf8]{inputenc}
\usepackage[T1]{fontenc}
\usepackage[french]{babel}
\usepackage{lmodern}
\usepackage[margin=2cm]{geometry}
\usepackage{graphicx}
\usepackage{xcolor}
\usepackage{booktabs}
\usepackage{listings}
\usepackage[hidelinks]{hyperref}

\lstset{basicstyle=\ttfamily\small, breaklines=true, frame=single, columns=fullflexible,
  literate={é}{{\'e}}1 {è}{{\`e}}1 {à}{{\`a}}1 {ê}{{\^e}}1 {ç}{{\c c}}1}

% Marque visible de ce qui reste a rediger
\newcommand{\afaire}[1]{\textcolor{red}{\textit{[À rédiger : #1]}}}

\begin{document}

% ------------------------------------------------------------------ En-tete
% Un simple bloc en haut de la page 1 (pas de page de titre : 3 pages maximum)
\begin{center}
  {\Large\bfseries Projet 1 --- Le robot de réconfort}\\[0.3cm]
  Master 1 Informatique --- POA --- projet de semestre\\[0.2cm]
  HAJEE Neha \quad --- \quad CAMARA ALY Zeinab\\[0.2cm]
  Enseignants : Vetea Jacot, Florian Lecourt --- Année universitaire 2026--2027\\[0.2cm]
  Dépôt : \url{https://github.com/nehahajee/POA_Projet1} --- Vidéo : \texttt{demo.mp4}
\end{center}

% ------------------------------------------------------------------
\section{Architecture}
% → enonce 11.1 : un schema, et la justification des frontieres entre
%   perception, memoire, decision et action.

\subsection{Schéma}
% → schema : Environnement (seul a lire la grille et les objets) <-> Robot
%   (perception / memoire / action) <- decision.py (decision), avec agent.py
%   qui charge, valide et assemble. affichage.py a part : il rejoue la trace.
% → \includegraphics[width=\linewidth]{schema_architecture.pdf}
\afaire{schéma}

\subsection{Rôle de chaque module}
% données :
%   agent.py             point d'entree : 4 arguments, validation, assemblage, ecriture de la trace
%   validation.py        controle du contenu des 4 fichiers (enonce 5.6), avant le lancement
%   environnement.py     le simulateur : seul a lire la grille et le contenu des casiers
%   robot.py             perception (percevoir), memoire (CarteMentale), actions (avancer,
%                        chercher, consulter, prendre, donner, attendre) et journalisation
%   deplacement.py       planification : BFS sur la carte mentale, rejoindre avec replanification
%   armoire.py           Roue de Plutchik, 15 candidats tries (6.2), trajet du selecteur, fouille
%   traitement_messages.py  normalisation du message et premier mot reconnu (6.1)
%   decision.py          enchainement d'une demande : dictionnaire, armoire, resident, PRENDRE, DONNER
\begin{tabular}{@{}lll@{}}
  \toprule
  Couche & Fichier / classe & Rôle \\
  \midrule
  Perception & \afaire{} & \afaire{} \\
  Mémoire    & \afaire{} & \afaire{} \\
  Décision   & \afaire{} & \afaire{} \\
  Action     & \afaire{} & \afaire{} \\
  \bottomrule
\end{tabular}

\subsection{Justification des frontières}
% → pourquoi le Robot ne recoit jamais la grille (connaissances = carte sans "grille",
%   enonce 3.3) et ne peut qu'interroger l'Environnement (percevoir, est_libre,
%   contenu_casier) ;
% → pourquoi la decision ne fait aucune action directement (elle passe par le Robot,
%   qui verifie la legalite et journalise : un pas illegal leve ValueError) ;
% → pourquoi la validation est separee de l'agent (elle lit la grille pour verifier le
%   FICHIER, mais ne dit jamais si un resident est atteignable) ;
% → ce que ca apporte pour le projet 2 (plusieurs robots partagent un Environnement).
\afaire{}

% ------------------------------------------------------------------
\section{Stratégie de replanification}
% → enonce 11.2 : que fait le robot quand il decouvre un mur ? Replanifie-t-il tout,
%   ou repare-t-il son plan ? Quel cout ?

\subsection{Planifier sur une carte incomplète}
% → hypothese optimiste : une case inconnue est supposee libre (annexe B.5) ;
%   BFS sur la carte mentale vers les cases libres adjacentes a la cible.
\afaire{}

\subsection{Ce qui se passe quand un mur est découvert}
% → apres chaque pas : perception, mise a jour de la memoire ; si la case suivante
%   du plan est un mur nouvellement vu, le robot s'arrete et RECALCULE tout le
%   chemin depuis sa position (on ne repare pas le plan) ;
% → difference avec l'annexe B.5, qui replanifie a CHAQUE pas : nous ne replanifions
%   que quand le plan est invalide ;
% données : appartement_01, trajet vers R1 : mur decouvert en [6, 7], detour par [4, 11]
%   (visible dans la video).
\afaire{}

\subsection{Coût}
% → nombre de BFS = 1 + nombre de murs qui coupent le plan ; chaque BFS coute
%   O(hauteur x largeur) ; termine car le nombre de murs est fini ;
% données : grand appartement de test 24 x 40 (cartes_test/test_grand.json) :
%   6 demandes, 1469 pas, 0,05 s d'execution ;
% données : impasse (cartes_test/test_impasse.json) : demi-tour, cout de l'hypothese optimiste.
\afaire{}

\subsection{Quand aucun chemin n'existe : la preuve d'inaccessibilité}
% → si le BFS echoue sur la carte mentale (inconnu = libre), aucun chemin reel
%   n'existe (annexe B.6) : le robot s'arrete et declare l'echec, sans jamais lire
%   la grille ;
% données : cartes_test/test_dico_inaccessible.json : preuve en 11 pas, puis 0 pas pour
%   la demande suivante (la memoire s'en souvient).
\afaire{}

% ------------------------------------------------------------------
\section{Critique de la règle de repli}
% → enonce 11.3 : la regle est-elle defendable ? En proposer une meilleure, et dire
%   ce qu'il faudrait changer dans les fichiers pour la mettre en oeuvre.

\subsection{Ce que fait la règle, sur un exemple}
% données : cartes_test/scenario_colonne_vide.json + donnees_test/armoire_tristesse_vide.json
%   (colonne tristesse vide) : « Je me sens tellement seule ce soir. » -> tristesse/moyenne,
%   repli voisine_1 sur [1, 5] = degout/moyenne -> objet « desodorisant » ;
%   le sens horaire departage surprise et degout, a egalite de distance.
\afaire{}

\subsection{Est-elle défendable ?}
% → arguments pour (simple, deterministe, verifiable par la trace) et contre (la
%   distance sur la roue n'est pas la proximite de reconfort : un objet "degout"
%   peut etre inadapte, voire blessant, pour une personne triste).
\afaire{}

\subsection{Une meilleure règle}
\afaire{proposition}

\subsection{Ce qu'il faudrait changer dans les fichiers}
% → armoire_standard.json et/ou dictionnaire.json : quels champs ajouter
%   (exemple : etiquettes multiples par objet, table de compatibilite entre
%   emotions) ; ce qui change dans armoire.py (calculer_candidats) et dans la trace.
\afaire{}

% ------------------------------------------------------------------
\section{Trois échecs de notre robot}
% → enonce 11.4 : trois situations ou il se comporte mal, avec l'analyse de la cause.
%   Pour chacun : la situation, la cause, ce que nous avons fait (ou ferions).

\subsection{Échec 1 : le résident muré bloquait tout le scénario}
% → situation : le robot prenait l'objet puis decouvrait que le resident etait mure ;
%   il ne peut pas reposer l'objet, donc plus de PRENDRE possible : demandes suivantes
%   bloquees.
% → correction : verifier l'accessibilite du resident AVANT PRENDRE, en marchant
%   (perception + memoire, jamais la grille), puis echec « resident inaccessible »
%   et passage a la demande suivante (choix valide par les enseignants).
% données : cout du detour sur scenario_01 : 160 pas sans verification, 260 avec ;
%   cartes_test/test_mure.json : demandes 2 et 3 en echec sans objet, demande 4 reussie.
\afaire{}

\subsection{Échec 2 : il retourne au dictionnaire pour un résident qu'il sait muré}
% → situation : 2e demande vers le meme resident mure ;
% → cause : la memoire n'est consultee qu'a l'etape « resident », apres le
%   dictionnaire et l'armoire ;
% données : cartes_test/test_mure.json, demande 3 : 53 pas pour un echec deja prouve ;
% → amelioration : tester d'abord le chemin sur la carte mentale (sans la grille).
\afaire{}

\subsection{Échec 3 : \afaire{titre}}
% → au choix, a partir de nos essais :
%   - l'impasse (cartes_test/test_impasse.json) : l'hypothese optimiste fait entrer le
%     robot dans un cul-de-sac avant de le voir ;
%   - la fouille de l'armoire repasse par des casiers deja vus vides (annexe C) ;
%   - le premier mot reconnu gagne : « je ne suis pas triste » donne tristesse.
\afaire{}

\end{document}
